\subsection{Continuous linear functions in a metrisable vector space}

THEOREM 1 (Banach). --- Let E and F be two metrisable vector spaces over a non-discrete valued division ring K, and let u be a continuous linear mapping of E in F. Suppose that E is complete. Then the following conditions are equivalent :

\begin{enumerate}
\def\labelenumi{(\roman{enumi})}
\item
  \(u\) is a strict surjective morphism.
\item
  F is complete and \(u\) is surjective.
\item
  The image of \(u\) is not meagre in \(\mathbf{F}\) (GT, IX, § 5.2).
\item
  For every neighbourhood \(\mathbf{V}\) of 0 in \(\mathbf{E}\) , the set \(\overline{u(\mathbf{V})}\) is a neighbourhood of 0 in \(\mathbf{F}\) . Firstly (i) implies (ii), for let \(u\) be a strict surjective morphism and \(\mathbf{N}\) be the kernel of \(u\) . Then \(u\) induces an isomorphism of \(\mathrm{E/N}\) on \(\mathbf{F}\) . But \(\mathbf{E}\) is metrisable and complete, hence \(\mathrm{E/N}\) is complete (GT, IX, § 3.1, prop. 4), therefore \(\mathbf{F}\) is complete.
\end{enumerate}

Next (ii) implies (iii). Let F be complete and u be surjective. The image of u is precisely F and therefore not meagre in F from Baire's theorem (GT, IX, § 5.3).

The following lemma shows that (iii) implies (iv).

Lemma 1. --- Let E and F be two topological vector spaces over a non-discrete valued division ring K, and let u be a continuous linear mapping of E in F such that the image of E is not meagre. Then, for every neighbourhood V of 0 in E, the set \(\overline{u(V)}\) is a neighbourhood of 0 in F.

Let W be a balanced neighbourhood of 0 in E such that \(W + W \subset V\) (I, § 1.5, prop. 4). Let \(\alpha\) be an element of K such that \(|\alpha| > 1\) ; then E is the union of the sets \(\alpha^{n}W\) where n varies in N; in fact, for all \(x \in E\) , there exists \(\beta \in K\) such that \(x \in \beta W\) (I, p.~7, prop. 4) and there exists an integer \(n \geqslant 0\) such that \(|\beta| < |\alpha|^{n}\) , then \(x \in \alpha^{n}W\) since W is balanced. Hence, \(u(E)\) is the union of the sequence of sets \(u(\alpha^{n}W) = \alpha^{n}u(W)\) , and as \(u(E)\) is not meagre in F, one at least of the sets \(\alpha^{n}\overline{u(W)}\) possesses an interior point (GT, IX, § 5.3, def. 2) and therefore \(\overline{u(W)}\) has an interior point.

Let \(y_0\) be an interior point of \(\overline{u(\mathbf{W})}\) ; since \(-u(\mathbf{W}) = u(\mathbf{W})\) , and therefore \(-\overline{u(\mathbf{W})} = \overline{u(\mathbf{W})}\) it follows that \(0 = y_0 + (-y_0)\) is an interior point of \(\overline{u(\mathbf{W})} + \overline{u(\mathbf{W})}\) . As vector addition is a continuous mapping of \(F \times F\) in \(F\) , the set \(\overline{u(\mathbf{W})} + \overline{u(\mathbf{W})}\) is contained in the closure of the set

\[
u (\mathrm{W}) + u (\mathrm{W}) = u (\mathrm{W} + \mathrm{W}) \subset u (\mathrm{V});
\]

hence \(\overline{u(\mathrm{V})}\) is a neighbourhood of 0 in F.

Before proving that (iv) implies (i) we prove the following lemma, where we make the convention that, in all metric spaces, \(\mathbf{B}_{r}(x)\) denotes the closed ball of centre x and radius r.

Lemma 2. --- Let E and F be two metric spaces, and suppose that E is also complete. Let u be a linear mapping of E in F having the following property: whatever the number r \textgreater{} 0, there exists a number \(\rho(r) > 0\) such that, for all \(x \in E\) , we have

\[
\mathrm{B} _ {\rho (r)} (u (x)) \subset \overline {{u (\mathrm{B} _ {r} (x))}}.
\]

In these conditions, for all \(a > r\) , the image \(u(\mathbf{B}_a(x))\) contains the ball \(\mathbf{B}_{\rho(r)}(u(x))\) .

Let \((r_n)\) be an infinite sequence of numbers \(>0\) such that \(r_1 = r\) and \(a = \sum_{n=1}^{\infty} r_n\) . For each index \(n\) there exists a number \(\rho_n > 0\) (with \(\rho_1 = \rho(r)\) ) such that

\[
\mathrm{B} _ {\rho_ {n}} (u (x)) \subset \overline {{u (\mathrm{B} _ {r _ {n}} (x))}}
\]

for all \(x \in E\) ; we can, and will, suppose that \(\lim \rho_{n} = 0\) .

Let \(x_0\) be a point of \(E\) , and \(y\) be a point of \(B_{\rho(r)}(u(x_0))\) . We shall show that \(y\) belongs to \(u(B_a(x_0))\) .

For this, a sequence \((x_{n})_{n>0}\) of points of E is defined inductively such that, for all \(n \geqslant 1\) , we have \(x_{n} \in \mathbf{B}_{r_{n}}(x_{n-1})\) and \(u(x_{n}) \in \mathbf{B}_{\rho_{n+1}}(y)\) . If the \(x_{i}\) have been defined for \(0 \leqslant i \leqslant n - 1\) satisfying these relations, then we have \(y \in \mathbf{B}_{\rho_{n}}(u(x_{n-1}))\) ; since

\[
\mathrm{B} _ {\rho_ {n}} (u (x _ {n - 1})) \subset \overline {{u (\mathrm{B} _ {r _ {n}} (x _ {n - 1}))}},
\]

there exists a point \(x_{n} \in \mathbf{B}_{r_n}(x_{n-1})\) whose image \(u(x_{n})\) belongs to the neighbourhood \(\mathbf{B}_{\rho_{n+1}}(y)\) of \(y\) , which establishes the existence of the sequence \((x_{n})\) .

Since the distance of \(x_{n}\) from \(x_{n+p}\) is less than \(r_{n+1} + r_{n+2} + \cdots + r_{n+p}\) , which is arbitrarily small when n is large, the sequence \((x_{n})\) is a Cauchy sequence in E. As E is complete, the sequence \((x_{n})\) converges to a point x of E. The distance of \(x_{0}\) from x is less than \(\sum_{n=1}^{\infty} r_{n} = a\) , thus \(x \in \mathbf{B}_{a}(x_{0})\) . But u is continuous, thus the sequence \(u(x_{n})\) converges to \(u(x)\) ; also \(u(x_{n}) \in \mathbf{B}_{\rho_{n+1}}(y)\) , hence \(y = u(x)\) , and the lemma is proved.

We return to the theorem and show that (iv) implies (i). Suppose that \(u\) satisfies condition (iv). For each of the spaces E and F, consider a distance that is invariant under translation and defines its topology (I, p.~16). By hypothesis, the set \(u(\mathbf{B}_r(0))\) is a neighbourhood of 0 in F for every \(r > 0\) , and thus there exists a number \(\rho(r) > 0\) such that \(\mathbf{B}_{\rho(r)}(0) \subset u(\mathbf{B}_r(0))\) . By translation we conclude that \(\mathbf{B}_{\rho(r)}(u(x)) \subset u(\mathbf{B}_r(x))\) for all \(r > 0\) and all \(x \in E\) . From lemma 2, for every pair of real positive numbers \((a, r)\) , \(a > r > 0\) , we have \(\mathbf{B}_{\rho(r)}(0) \subset u(\mathbf{B}_a(0))\) ; thus \(u\) is a strict morphism of E on F. We have shown that (iv) implies (i) and the proof of the theorem is completed.

COROLLARY 1. --- If E and F are two complete metrisable vector spaces over a non-discrete valued division ring, then every bijective continuous linear mapping of E on F is an isomorphism.

In particular, if E and F are two complete normed spaces, there exists a number \(a > 0\) such that \(\| u(x)\| \geqslant a.\| x\|\) for all \(x\in \mathbf{E}\) .

COROLLARY 2. --- Let E be a vector space over a non-discrete valued division ring, let \(T_{1}\) and \(T_{2}\) be two topologies on E compatible with its vector space structure and for each of which E is metrisable and complete. Then, if \(T_{1}\) and \(T_{2}\) are comparable, they are identical.

COROLLARY 3. --- Let E and F be two complete metrisable vector spaces over a non-discrete valued division ring. In order that a continuous linear mapping u of E in F should be a strict morphism, it is necessary and sufficient that \(u(\mathrm{E})\) be closed in F.

The condition is necessary, because if u is a strict morphism, the image \(u(\mathrm{E})\) , being isomorphic to the quotient \(\mathrm{E}/u^{-1}(0)\) , is complete (I, p.~17) and therefore closed in F. The condition is sufficient, since, if \(u(\mathrm{E})\) is closed in F, then \(u(\mathrm{E})\) must be a complete metrisable vector space and thus by theorem 1 u is a strict morphism of E on \(u(\mathrm{E})\) .

COROLLARY 4. --- Let E be a complete metrisable vector space over a non-discrete valued division ring. If M and N are two closed vector subspaces, that are (algebraic) complements in E, then E is the direct topological sum of M and N.

For \(M \times N\) is a complete metrisable vector space and the mapping \((y, z) \mapsto y + z\) of \(M \times N\) on E is continuous and bijective, therefore an isomorphism (cor. 1).

COROLLARY 5 (The closed graph theorem). --- Let E and F be two complete metrisable vector spaces over a non-discrete valued division ring. In order that a linear mapping u of E in F be continuous, it is necessary and sufficient that its graph, in the product space \(E \times F\) , be closed.

The condition is necessary since the graph of a continuous mapping into a Hausdorff space is closed (GT, I, § 8.1, cor. 2). To see that it is sufficient, note that it implies that the graph G of u, which is a closed vector subspace of the complete metrisable space \(E \times F\) , is itself metrisable and complete. The projection \(z \mapsto \operatorname{pr}_{1}(z)\) of G on E is a bijective, continuous linear mapping, therefore an isomorphism (cor. 1); since its inverse mapping is \(x \mapsto (x, u(x))\) , it follows that \(u\) is continuous in E.

We can express this corollary in the following form: \(u\) is continuous if the following situation holds: if the sequence \((x_{n})\) of points of E both converges to 0 and is such that the sequence \((u(x_{n}))\) converges to \(y\) , then it is necessarily the case that \(y = 0\) .

Example. --- Let E be a vector subspace of the space of real-valued functions defined on I = \([0, 1]\) ; let \(\|f\|\) be a norm on E, under which E is complete, and such that its topology is finer than the topology of simple convergence. Suppose further that E contains the set \(\mathcal{C}^{\infty}(\mathrm{I})\) of functions infinitely differentiable on I; we shall show that there exists an integer \(k \geqslant 0\) , such that E contains the set \(\mathcal{C}^{k}(\mathrm{I})\) of all functions with a continuous k-th derivative in I.

For every pair of integers \(m > 0\) , \(n \geqslant 0\) , let \(V_{mn}\) be the set of functions \(f \in \mathcal{C}^{\infty}(\mathrm{I})\) such that \(|f^{(h)}(x)| \leqslant 1 / m\) for \(0 \leqslant h \leqslant n\) and for all \(x \in \mathrm{I}\) . The \(V_{m,n}\) form a fundamental system of neighbourhoods of 0 for a metrisable topology compatible with the vector space structure of \(\mathcal{C}^{\infty}(\mathrm{I})\) , further \(\mathcal{C}^{\infty}(\mathrm{I})\) is complete in this topology (FVR, II, p.~2, th. 1). Let \(u\) be the canonical mapping of \(\mathcal{C}^{\infty}(\mathrm{I})\) in E; we show that \(u\) is continuous. From cor. 5 above it is sufficient to prove that if a sequence \((f_n)\) converges to 0 in \(\mathcal{C}^{\infty}(\mathrm{I})\) and to a limit \(f\) in E then necessarily \(f = 0\) . But this is immediate since, by hypothesis, \(f\) is the simple convergence limit of \((f_n)\) . Hence there exists an integer \(k \geqslant 0\) and a number \(a > 0\) such that the relation

\[
p_{k}(f) = \sup_{\substack{x\in \mathbf{I}\\ 0\leqslant h\leqslant k}}\big|f^{(h)}(x)\big|\leqslant a
\]

implies \(\| f\| \leqslant 1\) for all \(f\in \mathcal{C}^{\infty}(\mathrm{I})\)

But \(p_k\) is a norm on the space \(\mathcal{C}^k(\mathrm{I})\) and \(\mathcal{C}^\infty(\mathrm{I})\) is a subspace that is everywhere dense in \(\mathcal{C}^k(\mathrm{I})\) for this norm (the set of polynomials being already everywhere dense in \(\mathcal{C}^k(\mathrm{I})\) , an immediate consequence of the Weierstrass-Stone theorem). By what has gone before, the identity mapping of \(\mathcal{C}^\infty(\mathrm{I})\) (carrying the norm \(p_k\) ) in \(\mathrm{E}\) , is continuous, and so it can be extended continuously to the whole space \(\mathcal{C}^k(\mathrm{I})\) (since \(\mathrm{E}\) is complete). This proves our assertion.

PROPOSITION 1. --- Let E, F be two topological vector spaces over a non-discrete valued division ring K. We suppose that :

\begin{enumerate}
\def\labelenumi{\arabic{enumi})}
\item
  E is metrisable and complete.
\item
  There exists a sequence \((\mathbf{F}_{n})\) of complete metrisable vector spaces over K and, for each n, an injective continuous linear mapping \(v_{n}\) of \(F_{n}\) in F such that F is the union of the subspaces \(v_{n}(\mathbf{F}_{n})\) .
\end{enumerate}

Then let u be a linear mapping of E in F. If the graph of u is closed in \(E \times F\) , then there exists an integer n and a continuous linear mapping \(u_{n}\) of E in \(F_{n}\) such that \(u = v_{n} \circ u_{n}\) (which implies that u is continuous and that \(u(\mathrm{E}) \subset v_{n}(\mathrm{F}_{n})\) ).

Let G be the graph of u in \(E \times F\) . For all n, we consider the continuous linear mapping \(w_{n}:(x,y)\mapsto(x,v_{n}(y))\) of \(E \times F_{n}\) in \(E \times F\) ; as G is closed, the set \(w_{n}^{-1}(\mathbf{G})=\mathbf{G}_{n}\) is a closed vector subspace of \(E \times F_{n}\) ; if \(p_{n}\) is the restriction to \(G_{n}\) of the first projection \(pr_{1}\) , we have \(p_{n}(\mathbf{G}_{n})=u^{-1}(v_{n}(\mathbf{F}_{n}))\) . As \(p_{n}\) is continuous and \(G_{n}\) is complete (since \(G_{n}\) is closed in the complete space \(E \times F_{n}\) ), \(p_{n}(\mathbf{G}_{n})\) is, by theorem 1, either meagre in E or it is the whole of E. But, by hypothesis, E is the union of the \(p_{n}(\mathbf{G}_{n})\) , and as E is complete, the \(p_{n}(\mathbf{G}_{n})\) cannot all be meagre in E by Baire's theorem (GT, IX, § 5.3, th. 1). Therefore there exists an integer n such that \(p_{n}(\mathbf{G}_{n}) = \mathbf{E}\) , or in other words \(u(\mathbf{E}) \subset v_{n}(\mathbf{F}_{n})\) . Further, as \(v_{n}\) is injective, \(G_{n}\) is the graph of a linear mapping \(u_{n}\) of E in \(F_{n}\) , and by the closed graph theorem (I, p.~19, cor. 5) \(u_{n}\) is continuous; it follows then from the definitions that \(u = v_{n} \circ u_{n}\) .

